% ============================================================
%  Chapter 9 — GPS: Finding Your Way
%  Inside a Smartphone: The Tiny World in Your Hand
% ============================================================

\chapter{GPS: Finding Your Way}

\chapteropening{%
  You can be in the middle of a forest with no Wi-Fi and your phone
  still knows exactly where you are. How is that possible?
}
\chapterbadge{Mission: Navigate by Satellite!}

% --- Opening ---
\section*{Meet Orbi}

\character{Orbi}{Greetings from orbit! I am Orbi — one of dozens of GPS satellites
  circling the Earth right now, thousands of kilometres above you.
  I send a very precise signal to your phone every second.
  Your phone listens to me and a few of my friends, and together
  we tell it exactly where it is on Earth.}

\sectionrule

% --- What GPS is ---
\section*{What Is GPS?}

\term{GPS} stands for \textbf{Global Positioning System}.
It is a network of more than 30 satellites orbiting Earth.
They were originally built by the United States military,
but today anyone with a GPS receiver — including every smartphone — can use them for free.

\begin{funfact}
  GPS satellites orbit about \textbf{20,200 kilometres} above Earth.
  They travel so fast that they orbit the entire planet twice a day.
  Each satellite sends a signal that travels at the speed of light.
\end{funfact}

\sectionrule

% --- How it works ---
\section*{How Phones Know Where They Are}

Your phone has a tiny GPS \term{receiver} that listens for signals from satellites.
Each satellite constantly broadcasts two pieces of information:
\begin{itemize}
  \item Its exact location in space.
  \item The exact time the signal was sent.
\end{itemize}

The phone receives these signals and measures how long each signal took to arrive.
Since signals travel at the speed of light, even a tiny delay tells the phone
how far away each satellite is. The phone uses those distances to calculate its position.

\sectionrule

% --- Triangulation ---
\section*{Triangulation: Pinpointing Your Location}

The process of using multiple satellites to find a position is called \term{triangulation}.

\begin{picturethis}
  Imagine you are lost in a park and you can hear three different ice cream vans.
  You can tell which direction each van is and roughly how far away it sounds.
  Using all three distances, you can work out exactly where you must be standing.
  GPS works the same way — but with satellites instead of ice cream vans,
  and light-speed radio signals instead of music.
\end{picturethis}

\begin{quickmap}
  \textbf{How triangulation finds your location:}
  \begin{enumerate}[label=\textbf{\arabic*.}]
    \item Satellite A sends a signal. The phone knows it is \emph{some distance} from A.
    \item Satellite B sends a signal. The phone narrows the position down further.
    \item Satellite C sends a signal. The position is now pinpointed very precisely.
    \item A fourth satellite confirms the exact altitude (height above sea level).
  \end{enumerate}
\end{quickmap}

GPS gives your phone location numbers (latitude and longitude),
then your maps app turns those numbers into the map view you recognise.

\begin{diagram}[Three satellites triangulate a phone's position.]
  \begin{tikzpicture}[>=Stealth]
    % Earth (simplified)
    \fill[SkyBlue!40] (0,0) circle (1.4cm);
    \draw[SmartBlue,thick] (0,0) circle (1.4cm);
    \node[font=\small\bfseries] at (0,0) {Earth};
    % Satellites
    \node[draw=SunYellow!80!black,fill=SunYellow!20,circle,minimum size=0.6cm,
          font=\tiny\bfseries] (s1) at (90:4.2cm)  {Sat 1};
    \node[draw=SunYellow!80!black,fill=SunYellow!20,circle,minimum size=0.6cm,
          font=\tiny\bfseries] (s2) at (210:4.2cm) {Sat 2};
    \node[draw=SunYellow!80!black,fill=SunYellow!20,circle,minimum size=0.6cm,
          font=\tiny\bfseries] (s3) at (330:4.2cm) {Sat 3};
    % Dashed circles (range rings)
    \draw[dashed,SunYellow!60] (90:4.2cm)  circle (2.6cm);
    \draw[dashed,CoralRed!60]  (210:4.2cm) circle (2.6cm);
    \draw[dashed,NeonGreen]    (330:4.2cm) circle (2.6cm);
    % Phone at intersection
    \node[draw=TechPurple,fill=TechPurple!20,circle,minimum size=0.5cm,
          font=\tiny\bfseries,line width=1pt] at (0,-1.8) {Phone};
    % Signal lines
    \draw[SunYellow!80,->,thick] (s1.south) -- (0,-1.6);
    \draw[CoralRed!80,->,thick]  (s2.north east) -- (0,-1.6);
    \draw[NeonGreen,->,thick]    (s3.north west) -- (0,-1.6);
  \end{tikzpicture}
\end{diagram}

\sectionrule

\begin{newwords}
  \begin{description}[labelwidth=3.5cm, leftmargin=3.7cm]
    \item[\term{GPS}]              Global Positioning System — a satellite network for finding location.
    \item[\term{Satellite}]        A machine orbiting Earth that sends or receives signals.
    \item[\term{Receiver}]         A device that picks up GPS signals (built into every smartphone).
    \item[\term{Triangulation}]    Using signals from three or more sources to find an exact position.
  \end{description}
\end{newwords}

\sectionrule

\begin{tryit}
  \textbf{Find yourself on the map!}

  With a grown-up's help, open a maps app on a smartphone and allow it to use your location.
  \begin{itemize}
    \item Walk to a different room. Does the dot move?
    \item Go outside and walk half a block. How accurate is the position?
    \item Notice the small circle around the dot — that shows where the phone thinks you might be.
  \end{itemize}
  Can you figure out which direction is north by moving in a straight line?
\end{tryit}

\sectionrule

\begin{bigidea}
  \textbf{GPS} uses signals from multiple satellites orbiting Earth to calculate
  a phone's exact position through \textbf{triangulation}.
  It works anywhere on Earth, with no Wi-Fi or mobile signal needed.
\end{bigidea}

\character{Orbi}{You can now find your way anywhere on the planet!
  But with great power comes great responsibility.
  Our final chapter is about staying safe and smart. Follow Dot!}
